\section*{Hoofdstuk \ref{ch:g12:curly-arrows} --- Reactiemechanismen: gebogen pijlen}

\begin{solution}{exo:g12:curly-arrows:1}
\ce{OH-}: donor (\omterm{def:g10:lewis-and-shape:lone-pair}{vrije elektronenparen}, negatieve lading); geen
\omterm{def:g12:curly-arrows:site}{acceptorplaats}. \ce{H+}: acceptor (helemaal geen \omterm{def:g9:inside-the-atom:nucleus}{elektronen}).
\ce{CH2=CH2}: donor (de \omterm{def:g10:lewis-and-shape:multiple-bond}{dubbele binding}). \ce{CH3-Cl}: donor (\omterm{def:g10:lewis-and-shape:lone-pair}{vrije
elektronenparen} van \ce{Cl}), acceptor (het koolstofatoom,
$\delta^+$). \ce{H2O}: donor (\omterm{def:g10:lewis-and-shape:lone-pair}{vrije elektronenparen} van \ce{O}); zijn
waterstofatomen, $\delta^+$, zijn zwakke acceptorplaatsen.
\end{solution}

\begin{solution}{exo:g12:curly-arrows:2}
(a) \omterm{def:g12:curly-arrows:categories}{substitutie}; (b) \omterm{def:g12:curly-arrows:categories}{additie}; (c) \omterm{def:g12:curly-arrows:categories}{eliminatie}; (d) \omterm{def:g12:curly-arrows:categories}{additie}.
\end{solution}

\begin{solution}{exo:g12:curly-arrows:3}
Hij begint bij een elektronenpaar (een \omterm{def:g10:lewis-and-shape:lone-pair}{vrij elektronenpaar} of een
binding) en eindigt op het \omterm{def:g7:atoms-and-molecules:atom}{atoom} waar het paar naartoe gaat. Van een \omterm{def:g10:lewis-and-shape:lone-pair}{vrij
elektronenpaar} van zuurstof naar een koolstofatoom: het paar wordt een
nieuwe \ce{O-C}-binding.
\end{solution}

\begin{solution}{exo:g12:curly-arrows:4}
Ammoniak geeft het \omterm{def:g10:lewis-and-shape:lone-pair}{vrije elektronenpaar} van zijn stikstofatoom; de pijl
gaat van dat \omterm{def:g10:lewis-and-shape:lone-pair}{vrije elektronenpaar} naar \ce{H+}. Het product \ce{NH4+}
draagt de lading $+1$.
\end{solution}

\begin{solution}{exo:g12:curly-arrows:5}
Een deeltje dat in de ene stap ontstaat en in een latere stap verbruikt
wordt; het staat niet in de totale \omterm{def:g8:balanced-equations:equation}{reactievergelijking}. Hier het
carbokation \ce{CH3-CH2+}.
\end{solution}

\begin{solution}{exo:g12:curly-arrows:6}
Eén pijl van een \omterm{def:g10:lewis-and-shape:lone-pair}{vrij elektronenpaar} van het zuurstofatoom van \ce{OH-}
naar het koolstofatoom; één pijl van de \ce{C-Br}-binding naar het
broomatoom. Ladingen: $-1$ links
(\ce{OH-}); rechts is methanol neutraal en draagt \ce{Br-} $-1$:
in totaal $-1$ aan beide kanten.
\end{solution}

\begin{solution}{exo:g12:curly-arrows:7}
Broom is elektronegatiever dan waterstof: het neemt het paar mee en
vertrekt als \ce{Br-} ($-1$); het waterstofatoom bindt aan een
koolstofatoom, en het andere koolstofatoom houdt een positieve lading
over: \ce{CH3-CH2+} ($+1$).
\end{solution}

\begin{solution}{exo:g12:curly-arrows:8}
1.~\ce{CH3-CH2-OH + H+ -> CH3-CH2-OH2+}: donor, een \omterm{def:g10:lewis-and-shape:lone-pair}{vrij elektronenpaar}
van het zuurstofatoom; acceptor, \ce{H+}. 2.~\ce{CH3-CH2-OH2+ -> CH3-CH2+ + H2O}:
het \ce{C-O}-paar gaat naar het positieve zuurstofatoom (acceptor), dat
ermee vertrekt als water. 3.~\ce{CH3-CH2+ -> CH2=CH2 + H+}: het paar van
een \ce{C-H}-binding (donor) verschuift naar het positieve koolstofatoom
(acceptor) en vormt de \omterm{def:g10:lewis-and-shape:multiple-bond}{dubbele binding}, en \ce{H+} vertrekt. Het
waterstofion dat in stap 1 is gebruikt, wordt in stap 3 teruggegeven:
het staat aan beide kanten en valt weg.
\end{solution}

\begin{solution}{exo:g12:curly-arrows:9}
Een pijl van een \omterm{def:g10:lewis-and-shape:lone-pair}{vrij elektronenpaar} van het zuurstofatoom van water
naar het positieve koolstofatoom. Na het afsplitsen van \ce{H+}: ethanol
\ce{CH3-CH2-OH}. Totaal:
\ce{CH2=CH2 + H2O -> CH3-CH2-OH}, een \omterm{def:g12:curly-arrows:categories}{additie}.
\end{solution}

\begin{solution}{exo:g12:curly-arrows:10}
Er verschuiven twee paren, in één stap: de \ce{C-O}-binding ontstaat, de
\ce{C-Br}-binding wordt verbroken. Geen intermediair.
\end{solution}

\begin{solution}{exo:g12:curly-arrows:11}
Chloor is elektronegatiever ($3.16 - 2.55 = 0.61$ tegen
$2.96 - 2.55 = 0.41$): het koolstofatoom van chloormethaan is het meest
$\delta^+$, wat het op zichzelf sneller zou laten reageren. Dat doet het
niet: de snelheid hangt ook af van hoe gemakkelijk het \omterm{def:g10:electron-shells:family}{halogeen} met het
paar vertrekt, en het grotere broomatoom vertrekt gemakkelijker.
\end{solution}

\begin{solution}{exo:g12:curly-arrows:12}
Als het waterstofatoom aan koolstofatoom 1 bindt, ontstaat het
carbokation \ce{CH3-CH+-CH3}; dat geeft 2-chloorpropaan,
\ce{CH3-CHCl-CH3}. Als het aan koolstofatoom 2 bindt, ontstaat
\ce{+CH2-CH2-CH3}; dat geeft 1-chloorpropaan, \ce{CH2Cl-CH2-CH3}.
\end{solution}

\begin{solution}{exo:g12:curly-arrows:13}
De pijl wijst de verkeerde kant op: hij moet beginnen bij de donor, een
\omterm{def:g10:lewis-and-shape:lone-pair}{vrij elektronenpaar} van \ce{OH-}, en eindigen op het acceptor-koolstofatoom.
En de nieuwe binding zit tussen zuurstof en koolstof, niet tussen
zuurstof en broom: het broomatoom vertrekt en neemt het paar van de
\ce{C-Br}-binding mee.
\end{solution}

\begin{solution}{exo:g12:curly-arrows:14}
Het koolstofatoom van de \ce{C=O}-groep van een zuur is gebonden aan twee
elektronegatieve zuurstofatomen: sterk $\delta^+$, een acceptor. Het
zuurstofatoom van de \omterm{def:g11:functional-groups:alcohol}{alcohol} draagt $\delta^-$ en twee \omterm{def:g10:lewis-and-shape:lone-pair}{vrije
elektronenparen}: een donor. In totaal gaat er water af; de \ce{-OH} van
het zuur wordt vervangen door de \ce{-O-R} van de \omterm{def:g11:functional-groups:alcohol}{alcohol}: een
\omterm{def:g12:curly-arrows:categories}{substitutie}.
\end{solution}

\begin{solution}{exo:g12:curly-arrows:15}
Weggeduwd door de \omterm{def:g9:inside-the-atom:nucleus}{elektronen} van de \omterm{def:g10:lewis-and-shape:multiple-bond}{dubbele binding} verschuift het paar
van de \ce{Br-Br}-binding naar het verste broomatoom: het dichtstbijzijnde
broomatoom wordt $\delta^+$, een \omterm{def:g12:curly-arrows:site}{acceptorplaats}. Eerste stap: een pijl van
de \omterm{def:g10:lewis-and-shape:multiple-bond}{dubbele binding} naar het dichtstbijzijnde broomatoom, en een van de
\ce{Br-Br}-binding naar het verste broomatoom, dat vertrekt als
\ce{Br-}; op het andere koolstofatoom van de vroegere \omterm{def:g10:lewis-and-shape:multiple-bond}{dubbele binding}
blijft een positieve lading achter.
\end{solution}

\begin{solution}{pb:g12:curly-arrows:1}
\textbf{1.} Etheen: \ce{H2C=CH2}, een \omterm{def:g10:lewis-and-shape:multiple-bond}{dubbele binding} en vier \ce{C-H}.
Waterstofbromide: \ce{H-Br} met drie \omterm{def:g10:lewis-and-shape:lone-pair}{vrije elektronenparen} op het
broomatoom.

\textbf{2.} De \omterm{def:g10:lewis-and-shape:multiple-bond}{dubbele binding}: haar tweede paar is een gebied dat rijk
is aan \omterm{def:g9:inside-the-atom:nucleus}{elektronen}.

\textbf{3.} $2.96 - 2.20 = 0.76$: \ce{H} draagt $\delta^+$.

\textbf{4.} Het waterstofatoom.

\textbf{5.} Een pijl van de \omterm{def:g10:lewis-and-shape:multiple-bond}{dubbele binding} naar het waterstofatoom van
\ce{HBr}; een pijl van de \ce{H-Br}-binding naar het broomatoom.

\textbf{6.} Het carbokation \ce{CH3-CH2+} ($+1$) en \ce{Br-} ($-1$).

\textbf{7.} Ervoor: $0 + 0 = 0$; erna: $+1 - 1 = 0$.

\textbf{8.} Een pijl van een \omterm{def:g10:lewis-and-shape:lone-pair}{vrij elektronenpaar} van \ce{Br-} naar het
positieve koolstofatoom.

\textbf{9.} Zes (drie bindingen): één paar te weinig voor een octet, dus
het neemt een paar op van elke donor in de buurt.

\textbf{10.} Het carbokation \ce{CH3-CH2+}.

\textbf{11.} Het ontstaat in de eerste stap en wordt in de tweede
verbruikt.

\textbf{12.} Een \omterm{def:g12:curly-arrows:categories}{additie}.

\textbf{13.} Alleen de groep: de keten van twee koolstofatomen blijft
behouden, en de \omterm{def:g10:lewis-and-shape:multiple-bond}{dubbele binding} wordt een enkelvoudige binding met
\ce{H} en \ce{Br}.

\textbf{14.} \qty{28.0}{g/mol}; \qty{108.9}{g/mol}.

\textbf{15.} $10.0 / 28.0 = \qty{0.357}{mol}$.

\textbf{16.} Etheen: $0.357 - x$; \ce{HBr}: in overmaat; broomethaan: $x$.
Etheen is beperkend: $x_{\max} = \qty{0.357}{mol}$.

\textbf{17.} $0.357 \times 108.9 = \qty{38.9}{g}$.

\textbf{18.} $0.75 \times 38.9 = \qty{29.2}{g}$.

\textbf{19.} Ongeveer \qty{29}{g} broomethaan.
\end{solution}
